% Figure: the effect of a tuned mass damper. The steady-state response
% amplitude of a structure against forcing frequency, with and without
% an attached, tuned, damped absorber. The single resonance peak is
% split into two smaller peaks.
\begin{figure}[htbp]
\centering
\begin{tikzpicture}
\begin{axis}[
  width=12cm, height=7cm,
  xlabel={forcing frequency ratio $\omega/\omega_n$},
  ylabel={response amplitude (normalised)},
  xmin=0.5, xmax=1.5, ymin=0, ymax=10,
  axis lines=left,
  samples=300,
  legend pos=north east,
  legend cell align=left,
  grid=major, grid style={enggray!20},
  tick label style={font=\scriptsize},
  label style={font=\small},
  legend style={font=\scriptsize},
]
% bare structure, light damping zeta = 0.02: 1/sqrt((1-r^2)^2 + (2 zeta r)^2)
\addplot[engred, very thick, domain=0.5:1.5]
  {1/sqrt((1-x^2)^2 + (2*0.02*x)^2)};
\addlegendentry{bare structure}
% with tuned mass damper: a two-peaked curve, sketched with a smooth
% double-lobe. Use a fitted rational-like proxy that stays bounded.
\addplot[engblue, very thick, domain=0.5:1.5, samples=400]
  {2.6/sqrt((1 - (x/0.90)^2)^2 + (2*0.10*(x/0.90))^2)
   + 2.6/sqrt((1 - (x/1.11)^2)^2 + (2*0.10*(x/1.11))^2) - 1.4};
\addlegendentry{with tuned mass damper}
% mark the tuning frequency
\addplot[enggray, dashed, thick] coordinates {(1,0) (1,10)};
\node[font=\scriptsize, enggray, anchor=south west] at (axis cs:1.01,8.6)
  {tuned at $\omega_n$};
\end{axis}
\end{tikzpicture}
\caption[Response of a structure with a tuned mass damper]{The
steady-state response amplitude of a lightly damped structure driven
near its natural frequency $\omega_n$. Alone (red) the structure shows
one tall, sharp resonance peak that a matched forcing, wind gust or
footfall, can drive to damaging amplitude. Adding a smaller mass on a
spring and damper, tuned so its own natural frequency matches
$\omega_n$, splits the single peak into two lower ones (blue): the
absorber and the structure exchange energy out of phase, and the damper
on the absorber bleeds the vibration away. The absorber does not stiffen
the structure; it moves in opposition to it and carries the resonant
motion into its own damped mass. The device is why tall towers hang a
pendulum or a sliding block near the top, and why a footbridge that
sways underfoot is cured with a tuned absorber rather than a rebuild.}
\label{fig:vol03ch04:tmd-response}
\end{figure}
